\section{Data sources and descriptives}

The dataset is built from Spain's Labour Force Survey microdata (EPA), published by INE. Quarterly microdata files are downloaded for 2006--2025 and processed into annual estimates by age group and birthplace status for individuals aged 16 and over. Individuals are classified as Spain-born when their reported place of birth corresponds to a Spanish province, and as foreign-born otherwise. Population totals are computed using EPA survey weights. Labour-force participation is measured as the weighted share of individuals who are employed or unemployed, while working hours are measured as weighted usual weekly hours among employed workers. Annual labour input is then calculated by combining population, participation, and weekly hours, assuming a 50-week working year.

\vspace{0.7em}\par
The original EPA microdata report age in grouped intervals—16--19, 20--24, and successive five‑year bands up to 65+. These groups are kept when computing annual estimates and decomposition effects. However in summary tables and plots, labour supply and contributions to change, estimates are aggregated into three broader categories: 16--34 (young), 35--54 (adult), and 55+ (older).The annex contain description of the above variable related to survey quetionaire.


\vspace{0.7em}\par
Table~\ref{tab:popSum} reports summary indicators by birthplace status for Spain in 2005 and 2025, including population, activity rates, average weekly hours, and the share of the population aged 55 and over. The estimates are calculated from EPA microdata using survey weights, so they represent weighted population estimates rather than official population totals. They may therefore differ from official residency statistics because of differences in definitions and methodology. One important distinction is that official residency estimates are commonly reported by nationality, whereas this analysis classifies individuals by place of birth.

\begin{table}[H]%
  \caption{Population aged 16 and over, activity, hours, and ageing indicators by birthplace, Spain, 2005 and 2025.}

  %\source{Author's calculations}
  \label{tab:popSum}%
  \begin{center}
    \begin{tabular*}{\textwidth}{@{\extracolsep{\fill}}lcrrr}
\toprule
\multicolumn{1}{c}{\bfseries }&\multicolumn{1}{c}{\bfseries }&\multicolumn{3}{c}{\bfseries Birthplace status}\tabularnewline
\cline{3-5}
\multicolumn{1}{l}{}&\multicolumn{1}{c}{}&\multicolumn{1}{r}{Spanish}&\multicolumn{1}{r}{Foreign}&\multicolumn{1}{r}{Combined}\tabularnewline
\midrule
\textbf{Population}&&&&\tabularnewline
\quad 2005&&32.79&3.79&36.58\tabularnewline
\quad 2025&&33.58&8.53&42.11\tabularnewline
\quad Change&&0.79&4.74&5.53\tabularnewline
\textbf{Activity rate}&&&&\tabularnewline
\quad 2005&&55.83&74.77&57.79\tabularnewline
\quad 2025&&55.79&71.44&58.96\tabularnewline
\quad Change&&-0.04&-3.33&1.17\tabularnewline
\textbf{Weekly hours}&&&&\tabularnewline
\quad 2005&&42.63&42.69&42.64\tabularnewline
\quad 2025&&39.16&40.11&39.35\tabularnewline
\quad Change&&-3.47&-2.58&-3.29\tabularnewline
\textbf{Share aged 55+}&&&&\tabularnewline
\quad 2005&&34.11&11.01&31.71\tabularnewline
\quad 2025&&45.02&23.85&40.73\tabularnewline
\quad Change&&10.91&12.84&9.02\tabularnewline
\bottomrule
\end{tabular*}
%
    \caption*{\small Note: Population is in millions; activity rate and the share aged 55+ are percentages; weekly hours are average hours worked per week; Combined ativity rate and weekly hours are computed as population‑weighted means of the age‑group measures.}
  \end{center}
\end{table}%

Table~\ref{tab:popSum} shows that foreign-born individuals became a much larger part of Spain's labour supply between 2005 and 2025. Their population more than doubled, from 3.79 to 8.53 million, while the Spain-born population changed little. Foreign-born individuals also maintained higher activity rates than Spain-born individuals, although their average weekly hours declined slightly. Both groups aged over the period, with the share aged 55+ rising especially among the foreign-born population. As a result, the foreign-born share of the total population increased from 10.36\% in 2005 to 20.26\% in 2025.



